\chapter{The Yoga Texts}
\label{ch:texts}

\emph{Which texts does the syllabus name, and what must you know from each?}
Learn the named sutras and slokas with their sense, the three foundation texts of Vedanta, and how the texts speak of health, the mind and the eight limbs.
Keep the Sanskrit term with one English word. The Devanagari below is still to be set against a printed edition before print, and a certified teacher must review this chapter.

\section{Patanjala Yoga Sutra}
\label{sec:patanjala-yoga-sutra}

The Yoga Sutra is a compact collection of aphorisms in four chapters (padas):
\begin{itemize}
\item samadhi, on absorption;
\item sadhana, on practice;
\item vibhuti, on the powers that arise;
\item kaivalya, on freedom.
\end{itemize}
The syllabus names the first twelve sutras of the first pada. Learn each with its sense.

\subsection*{The twelve sutras with their sense}
\begin{itemize}
\item atha yogānuśāsanam: now begins the instruction on yoga.
\item yogaś citta-vṛtti-nirodhaḥ: yoga is the stilling of the movements of the mind.
\item tadā draṣṭuḥ svarūpe 'vasthānam: then the seer rests in its own nature.
\item vṛtti-sārūpyam itaratra: at other times the seer seems to take the shape of the mind's movements.
\item vṛttayaḥ pañcatayyaḥ kliṣṭākliṣṭāḥ: the movements are of five kinds, some painful and some not.
\item pramāṇa-viparyaya-vikalpa-nidrā-smṛtayaḥ: the five are right knowledge, error, imagination, sleep and memory.
\item pratyakṣānumānāgamāḥ pramāṇāni: right knowledge is perception, inference and reliable testimony.
\item viparyayo mithyā-jñānam atad-rūpa-pratiṣṭham: error is false knowledge that does not rest on the real form of a thing.
\item śabda-jñānānupātī vastu-śūnyo vikalpaḥ: imagination is a notion that follows words but has no object.
\item abhāva-pratyayālambanā vṛttir nidrā: sleep is the movement that rests on the sense of absence.
\item anubhūta-viṣayāsampramoṣaḥ smṛtiḥ: memory is not letting go of an object once experienced.
\item abhyāsa-vairāgyābhyāṃ tan-nirodhaḥ: their stilling comes through practice and dispassion.
\end{itemize}

\subsection*{Text for recitation}
\begin{itemize}
\item[1.1] \sa{अथ योगानुशासनम्} (atha yogānuśāsanam)
\item[1.2] \sa{योगश्चित्तवृत्तिनिरोधः} (yogaś cittavṛttinirodhaḥ)
\item[1.3] \sa{तदा द्रष्टुः स्वरूपेऽवस्थानम्} (tadā draṣṭuḥ svarūpe 'vasthānam)
\item[1.4] \sa{वृत्तिसारूप्यमितरत्र} (vṛttisārūpyam itaratra)
\item[1.5] \sa{वृत्तयः पञ्चतय्यः क्लिष्टाक्लिष्टाः} (vṛttayaḥ pañcatayyaḥ kliṣṭākliṣṭāḥ)
\item[1.6] \sa{प्रमाणविपर्ययविकल्पनिद्रास्मृतयः} (pramāṇaviparyayavikalpanidrāsmṛtayaḥ)
\item[1.7] \sa{प्रत्यक्षानुमानागमाः प्रमाणानि} (pratyakṣānumānāgamāḥ pramāṇāni)
\item[1.8] \sa{विपर्ययो मिथ्याज्ञानमतद्रूपप्रतिष्ठम्} (viparyayo mithyājñānam atadrūpapratiṣṭham)
\item[1.9] \sa{शब्दज्ञानानुपाती वस्तुशून्यो विकल्पः} (śabdajñānānupātī vastuśūnyo vikalpaḥ)
\item[1.10] \sa{अभावप्रत्ययालम्बना वृत्तिर्निद्रा} (abhāvapratyayālambanā vṛttir nidrā)
\item[1.11] \sa{अनुभूतविषयासंप्रमोषः स्मृतिः} (anubhūtaviṣayāsaṃpramoṣaḥ smṛtiḥ)
\item[1.12] \sa{अभ्यासवैराग्याभ्यां तन्निरोधः} (abhyāsavairāgyābhyāṃ tannirodhaḥ)
\end{itemize}

\subsection*{How the twelve fit together}
They move from a statement of the aim, through the way the mind moves, to the two means of stilling it. Teach them in that order: the aim (1.1 to 1.4), the way the mind moves (1.5 to 1.11), the two means (1.12). Have each student say one sutra aloud, then its sense in plain words.

\subsection*{Two more points from the Sutra}
\begin{itemize}
\item Five kleshas (2.3): avidya (ignorance), asmita (ego-sense), raga (liking), dvesha (dislike), abhinivesha (clinging to life).
\item Posture (2.46): sthira-sukham asanam, a steady and comfortable posture.
\end{itemize}

\begin{recall}
\begin{itemize}
\item Four padas: samadhi, sadhana, vibhuti, kaivalya.
\item 1.2: yoga is the stilling of the movements of the mind.
\item Five movements (1.6): right knowledge, error, imagination, sleep, memory.
\item Right knowledge (1.7): perception, inference, reliable testimony.
\item Two means (1.12): practice (abhyasa) and dispassion (vairagya).
\item Five kleshas (2.3): ignorance, ego-sense, liking, dislike, clinging to life.
\end{itemize}
\end{recall}

\subsection*{Common questions}
\begin{itemize}
\item \emph{Which sutras must I know?} The first twelve of the first pada, with their sense.
\item \emph{Is the Devanagari final?} Not yet. It was converted mechanically from the checked romanised text, so a printed edition should confirm the spelling and the avagraha in 1.3 before print.
\item \emph{Is the Sanskrit in this chapter reliable?} The sutra text and numbers were checked against the GRETIL edition of the Vyasa bhashya text. Set it against a printed edition before print.
\end{itemize}

\subsection*{Worked examples}
\begin{itemize}
\item \textbf{Model answer: state and explain Patanjali's definition of yoga.} Quote the sutra with its number: yogaś citta-vṛtti-nirodhaḥ (1.2). Give the meaning: yoga is the stilling of the movements of the mind. Link to 1.3: then the seer rests in its own nature. Link to 1.4: at other times the seer seems to take the shape of the mind's movements.
\item \textbf{Memory aid for the five movements.} Say the sutra pramāṇa-viparyaya-vikalpa-nidrā-smṛtayaḥ, which lists the five in order. Then give one line for each from the next sutras: 1.7 the three sources of right knowledge; 1.8 error is false knowledge; 1.9 imagination follows words but has no object; 1.10 sleep rests on the sense of absence; 1.11 memory is not letting go of an object once experienced.
\item \textbf{Reciting the twelve in order.} Follow the order above, aim, then movements of the mind, then the two means. Each student says one sutra aloud, then its sense, and moves on in sequence.
\end{itemize}

\subsection*{Check yourself}
\begin{enumerate}
\item Name the four padas of the Yoga Sutra.
\item Fill in the blank: the five movements of the mind are right knowledge, error, \underline{\hspace{2cm}}, sleep and memory.
\item Which two means stop the movements of the mind, and in which sutra?
\item Name the five kleshas.
\item Match the sutra with its sense: 1.1, 1.2, 1.3. Senses: yoga is the stilling of the movements of the mind; now begins the instruction on yoga; then the seer rests in its own nature.
\item True or false: 1.7 says right knowledge comes from perception and memory.
\item What phrase in 2.46 describes the posture, and what does it mean?
\item A student says error and imagination are the same, both are wrong ideas. Correct this using 1.8 and 1.9.
\item Explain in your own words how the twelve sutras fit together.
\item How many kinds of mental movement does the Sutra name, and which are they?
\end{enumerate}

\subsection*{Answers}
\begin{enumerate}
\item Samadhi (absorption), sadhana (practice), vibhuti (the powers that arise) and kaivalya (freedom).
\item Imagination (vikalpa).
\item Practice (abhyasa) and dispassion (vairagya), in 1.12.
\item Avidya (ignorance), asmita (ego-sense), raga (liking), dvesha (dislike) and abhinivesha (clinging to life), from 2.3.
\item 1.1 now begins the instruction on yoga; 1.2 yoga is the stilling of the movements of the mind; 1.3 then the seer rests in its own nature.
\item False. Right knowledge is perception, inference and reliable testimony. Memory is a separate movement of the mind.
\item Sthira-sukham asanam, a steady and comfortable posture.
\item Error is false knowledge that does not rest on the real form of a thing. Imagination is a notion that follows words but has no object. They are two separate movements.
\item They begin with the statement of the aim, go through the way the mind moves, and end with the two means of stilling it, practice and dispassion.
\item Five: right knowledge, error, imagination, sleep and memory.
\end{enumerate}

\paragraph{Sources.} The Board syllabus for Level 1 and the official guidebook for Level I (2016, older scheme, coverage check only). The sutra text was read in the GRETIL edition of the Vyasa bhashya text; a printed edition with a standard commentary is still to be set against it.

\section{Bhagavad Gita}
\label{sec:bhagavad-gita}

The Gita is a dialogue between Krishna and Arjuna in the Mahabharata. It is read as a guide to action, devotion and knowledge, and it describes yoga as even-mindedness and skill in action.
The syllabus names five slokas of the second chapter: 47, 48, 49, 50 and 70.

\subsection*{The five slokas with their sense}
\begin{itemize}
\item 47: you have a right to the action, never to its fruits; do not act for the sake of fruit, and do not cling to inaction.
\item 48: act established in yoga, giving up attachment and staying even in success and failure; evenness is called yoga.
\item 49: action driven by desire for results is far lower than the yoga of a steady mind; take refuge in the steady mind, for those who act for fruit are to be pitied.
\item 50: one joined to the steady mind lets go of both good and bad results; so strive for yoga, which is skill in action.
\item 70: as waters flow into an ocean that is full and still, so desires flow into the one who stays at peace; the one who runs after desires does not find peace.
\end{itemize}

\subsection*{Text for recitation}
\begin{itemize}
\item[2.47] \sa{कर्मण्येवाधिकारस्ते मा फलेषु कदाचन । मा कर्मफलहेतुर्भूर्मा ते सङ्गोऽस्त्वकर्मणि ॥} (karmaṇy-evādhikāras te mā phaleṣu kadācana / mā karma-phala-hetur bhūr mā te saṅgo 'stv akarmaṇi)
\item[2.48] \sa{योगस्थः कुरु कर्माणि सङ्गं त्यक्त्वा धनञ्जय । सिद्ध्यसिद्ध्योः समो भूत्वा समत्वं योग उच्यते ॥} (yoga-sthaḥ kuru karmāṇi saṅgaṁ tyaktvā dhanañjaya / siddhy-asiddhyoḥ samo bhūtvā samatvaṁ yoga ucyate)
\item[2.49] \sa{दूरेण ह्यवरं कर्म बुद्धियोगाद्धनञ्जय । बुद्धौ शरणमन्विच्छ कृपणाः फलहेतवः ॥} (dūreṇa hy avaraṁ karma buddhi-yogād dhanañjaya / buddhau śaraṇam anviccha kṛpaṇāḥ phala-hetavaḥ)
\item[2.50] \sa{बुद्धियुक्तो जहातीह उभे सुकृतदुष्कृते । तस्माद्योगाय युज्यस्व योगः कर्मसु कौशलम् ॥} (buddhi-yukto jahātīha ubhe sukṛta-duṣkṛte / tasmād yogāya yujyasva yogaḥ karmasu kauśalam)
\item[2.70] \sa{आपूर्यमाणमचलप्रतिष्ठं समुद्रमापः प्रविशन्ति यद्वत् । तद्वत्कामा यं प्रविशन्ति सर्वे स शान्तिमाप्नोति न कामकामी ॥} (āpūryamāṇam acala-pratiṣṭhaṁ samudram āpaḥ praviśanti yadvat / tadvat kāmā yaṁ praviśanti sarve sa śāntim āpnoti na kāma-kāmī)
\end{itemize}

\subsection*{The one idea}
The five slokas build one idea: act fully, drop the grip on results, stay even, and find peace. Give one workplace or classroom example for each.

\begin{recall}
\begin{itemize}
\item Krishna and Arjuna, in the Mahabharata.
\item Slokas to know: 2.47, 2.48, 2.49, 2.50, 2.70.
\item Two definitions of yoga: evenness of mind in success and failure (2.48), and skill in action (2.50, karmasu kauśalam).
\item Act, but do not be driven by the fruit; do not cling to inaction.
\end{itemize}
\end{recall}

\subsection*{Common questions}
\begin{itemize}
\item \emph{Which verses must I recite?} 2.47, 2.48, 2.49, 2.50 and 2.70. Learn the text with its meaning and give the sense in your own words.
\item \emph{Is the recitation text reliable?} The Devanagari, transliteration and sense of every verse were read against one edition. A second edition is still to be set against it before print.
\item \emph{The workshop slides show different text for 49, 50 and 70. Which do I follow?} Follow the recitation text here. The slide text for those verses was not checked.
\item \emph{Which text defines yoga as skill in action?} The Bhagavad Gita, in 2.50.
\end{itemize}

\subsection*{Worked examples}
\begin{itemize}
\item \textbf{Model answer: explain Gita 2.47.} The verse is karmaṇy-evādhikāras te mā phaleṣu kadācana. You have a right to the action, never to its fruits. Do not act for the sake of fruit. Do not cling to inaction. Link: 2.48 says evenness in success and failure is called yoga.
\item \textbf{Memory aid for the five slokas.} Act fully (47). Drop the grip on results and stay even (48). Prefer the steady mind to action driven by desire (49). Let go of both good and bad results, and see yoga as skill in action (50). Find peace, like the ocean that stays still as waters flow in (70). Say the five in order with one line each.
\item \textbf{Two definitions of yoga, one answer.} Asked how the Gita defines yoga, answer with two verses: 2.48, evenness of mind (samatvam) in success and failure, and 2.50, yoga is skill in action (karmasu kauśalam). Give the verse number with each.
\end{itemize}

\subsection*{Check yourself}
\begin{enumerate}
\item Who speak in the Gita, and in which larger work does it appear?
\item Which five slokas of the second chapter does the syllabus name?
\item Fill in the blank: in 2.50, yoga is \underline{\hspace{2cm}} in action.
\item True or false: 2.47 tells us to give up action.
\item Match the verse with its main idea: 2.48, 2.49, 2.70. Ideas: evenness is yoga; desires flow into the one at peace, like waters into a full still ocean; action driven by desire for results is far lower than the yoga of a steady mind.
\item What image does 2.70 use, and what does it say about the one who runs after desires?
\item A student works hard for exams but is always worried about the result. Which verse speaks to this, and what does it advise?
\item A student says the Gita teaches do nothing and accept what comes. Correct this.
\item What are the two definitions of yoga in the second chapter?
\end{enumerate}

\subsection*{Answers}
\begin{enumerate}
\item Krishna and Arjuna, in the Mahabharata.
\item 47, 48, 49, 50 and 70.
\item Skill (karmasu kauśalam).
\item False. It says we have a right to the action, never to its fruits, and we must not cling to inaction either.
\item 2.48 evenness is yoga; 2.49 action driven by desire for results is far lower than the yoga of a steady mind; 2.70 desires flow into the one at peace, like waters into a full still ocean.
\item Waters flowing into an ocean that is full and still. The one who runs after desires does not find peace.
\item 2.47: act, but do not be driven by the fruit. Together with 2.48, stay even in success and failure.
\item 2.47 says not to cling to inaction, and 2.50 calls yoga skill in action. The Gita asks for full action without grip on the results.
\item Evenness of mind in success and failure, and skill in action.
\end{enumerate}

\paragraph{Sources.} The Board syllabus for Level 1. The verse text, transliteration and sense were read against one edition (the holy-bhagavad-gita.org text of Swami Mukundananda). The official guidebook has no section on the Gita, so a second standard edition with commentary is still to be set against this text.

\section{Prasthanatrayi}
\label{sec:prasthanatrayi}

The Prasthanatrayi is the set of three sources (prasthanas, starting points) on which the Vedanta tradition rests. Remember the three by name and by one word each: heard, remembered, reasoned.

\subsection*{The three}
\begin{itemize}
\item The principal Upanishads: the shruti prasthana, what is heard, the revealed text.
\item The Bhagavad Gita: the smriti prasthana, what is remembered.
\item The Brahma Sutras, also called the Vedanta Sutras: the nyaya prasthana, also tarka prasthana, the reasoned and logical statement of the Upanishadic teaching.
\end{itemize}

\subsection*{Why it matters for yoga}
The Gita and the Upanishads carry much of the teaching on yoga. The three are the base on which the schools of Vedanta build.

\subsection*{A common teaching device}
The first six chapters of the Gita are taken as karma yoga, the next six as bhakti yoga and the last six as jnana yoga. This is a teaching convention, not the text's own division.

\begin{recall}
\begin{itemize}
\item Prasthana means a source or starting point.
\item Upanishads: shruti, heard. Gita: smriti, remembered. Brahma Sutras: nyaya or tarka, reasoned.
\item Brahma Sutras and Vedanta Sutras are the same text.
\item The six-six-six grouping of the Gita is a convention.
\end{itemize}
\end{recall}

\subsection*{Common questions}
\begin{itemize}
\item \emph{Do I need the list of principal Upanishads or verse counts?} No. The syllabus does not ask for them, and the number of principal Upanishads varies between lists.
\item \emph{Does the official guidebook cover this?} No. The three texts and their names are stated the same way in standard textbooks; the Brahma Sutras are attributed to Badarayana, in four chapters of four parts each.
\end{itemize}

\subsection*{Worked examples}
\begin{itemize}
\item \textbf{Model answer: what is the Prasthanatrayi?} Give the meaning first: the three prasthanas on which the Vedanta tradition rests. Then the three: the principal Upanishads, the Bhagavad Gita and the Brahma Sutras. Then the role of each: shruti, smriti, nyaya or tarka. Close with why it matters: the three are the base on which the schools of Vedanta build.
\item \textbf{Memory aid: heard, remembered, reasoned.} Heard is what was received, so the Upanishads. Remembered is smriti, so the Gita. Reasoned is the logical statement of the Upanishadic teaching, so the Brahma Sutras. If the exam gives a prasthana name, turn it into the one word first and the text follows.
\item \textbf{A student asks why the Gita has three yogas.} Say the Gita is one of the three, the smriti prasthana. Then explain the six-six-six teaching device, and add that it is a convention and not the text's own division.
\end{itemize}

\subsection*{Check yourself}
\begin{enumerate}
\item Fill in the blank: the Prasthanatrayi is made of the principal Upanishads, the Bhagavad Gita and the \underline{\hspace{2cm}}.
\item Which text is the shruti prasthana?
\item Match each text to its prasthana: Upanishads, Gita, Brahma Sutras with nyaya, shruti, smriti.
\item True or false: the Gita is the nyaya prasthana.
\item What does nyaya or tarka prasthana tell you about the Brahma Sutras?
\item How is the Gita grouped by the common teaching device?
\item A student says the six-six-six grouping is how the Gita itself divides its chapters. How do you correct this?
\item In two lines, explain why the Prasthanatrayi matters for a yoga candidate.
\end{enumerate}

\subsection*{Answers}
\begin{enumerate}
\item The Brahma Sutras (also called the Vedanta Sutras).
\item The Upanishads (what is heard, the revealed text).
\item Upanishads with shruti; Gita with smriti; Brahma Sutras with nyaya (also called tarka).
\item False. The Gita is the smriti prasthana, what is remembered. The Brahma Sutras are the nyaya prasthana.
\item They are the reasoned and logical statement of the Upanishadic teaching.
\item The first six chapters as karma yoga, the next six as bhakti yoga, the last six as jnana yoga.
\item Say that it is a teaching convention, a common device, and not the text's own division.
\item The Gita and the Upanishads carry much of the teaching on yoga, and the three texts together are the base on which the schools of Vedanta build.
\end{enumerate}

\paragraph{Sources.} The Board syllabus for Level 1; standard references on the Prasthanatrayi and the Brahma Sutras (attributed to Badarayana; four chapters of four parts each). The six-six-six grouping of the Gita is a teaching convention.

\section{How hatha practices support health}
\label{sec:hatha-practices-health}

Hatha practices are meant to balance the body and mind and to prepare for meditation, not only for treatment. They are mainly preventive. Say prevention and well-being, never cure.

\subsection*{What each practice adds}
This follows the Hatha Pradipika (second chapter) and the Gheranda Samhita; verse numbers differ between editions, so the chapter is cited rather than the verse.
\begin{itemize}
\item Kriyas give purification.
\item Asanas give strength and health (arogya).
\item Pranayama gives lightness.
\item Mudras and bandhas give steadiness.
\end{itemize}

\subsection*{Three levels of benefit}
\begin{itemize}
\item \textbf{Body:} flexibility, strength, balance and better alignment from asana; healthier breathing from pranayama; cleansing of throat, stomach, bowels and lungs from the shatkarma.
\item \textbf{Mind:} attention on the practice takes the mind off worries; pranayama brings calmness; practice makes one more aware of habitual thoughts and reactions.
\item \textbf{Inner side:} an easier meditation, inner peace, and growth towards self-knowledge, by joining physical and inner practice.
\end{itemize}

\subsection*{Everyday benefits of asana}
Calmness, flexibility, strength, balance, better body awareness, and less worry while attention is on the practice.

\subsection*{Weight and fitness}
Yoga alone burns fewer calories than strong exercise. With sensible diet and some aerobic activity it can help control weight.

\subsection*{Middle path}
Neither fasting nor overeating, neither indulgence nor deprivation, no wasting of energy.

\subsection*{General rules for asana practice}
\begin{itemize}
\item Empty or light bowel and bladder.
\item A ventilated place.
\item A level floor with a mat or blanket.
\item Unwanted thoughts kept aside.
\item No practice during illness, after surgery, fracture or sprain until a physician allows.
\end{itemize}

\begin{caution}
Do not claim that hatha yoga treats or cures any condition, reduces anxiety, detoxifies the body or improves stamina and circulation. The guidebook does not state stamina and circulation, and detoxification is left out as a health claim. A certified teacher or qualified reviewer must review this section.
\end{caution}

\begin{recall}
\begin{itemize}
\item Kriya: purification. Asana: strength and health. Pranayama: lightness. Mudra and bandha: steadiness. Say it as a chain: clean, strong, light, steady.
\item Three levels: body, mind, inner side.
\item Mainly preventive; prepares for meditation.
\item Meditation is a goal that hatha prepares for. The four limbs of hatha are asana, pranayama, mudra and bandha, and nada-anusandhana.
\end{itemize}
\end{recall}

\subsection*{Common questions}
\begin{itemize}
\item \emph{Is hatha yoga a treatment?} No. It is mainly preventive and aims to balance body and mind and prepare for meditation.
\item \emph{Does it help detoxify the body?} The guidebook names cleansing of throat, stomach, bowels and lungs from the shatkarma, but detoxification is left out as a health claim.
\item \emph{Is meditation part of hatha yoga?} No. Meditation is a goal that hatha prepares for.
\end{itemize}

\subsection*{Worked examples}
\begin{itemize}
\item \textbf{Model answer: significance of hatha practices in health and well-being.} Aim: balance body and mind and prepare for meditation, mainly preventive. Body: flexibility, strength, balance, alignment, healthier breathing, cleansing from the shatkarma. Mind: attention takes the mind off worries, pranayama brings calmness, more awareness of habitual thoughts. Inner side: easier meditation, inner peace, growth towards self-knowledge. Do not claim cure.
\item \textbf{A student with a sprained ankle wants to practise at once.} Recall the rule: no practice during illness, after surgery, fracture or sprain until a physician allows. Ask the student to get the physician's approval first. Remind them the aim is balance and prevention, so waiting does not lose it.
\end{itemize}

\subsection*{Check yourself}
\begin{enumerate}
\item Fill in the blank: kriya gives purification, asana gives strength and health, pranayama gives lightness, mudra and bandha give \underline{\hspace{2cm}}.
\item What are the three levels of benefit of hatha yoga?
\item True or false: hatha yoga is mainly meant to treat disease.
\item Give three everyday benefits of asana.
\item Match: (a) kriya, (b) pranayama, (c) asana, (d) mudra and bandha, with (1) lightness, (2) steadiness, (3) purification, (4) strength and health.
\item Can yoga alone control weight?
\item Name four general rules to follow before asana practice.
\item A student asks if hatha yoga reduces anxiety. How should a teacher word the answer?
\end{enumerate}

\subsection*{Answers}
\begin{enumerate}
\item Steadiness.
\item Body, mind and the inner (spiritual) side.
\item False. It is mainly preventive and prepares for meditation.
\item Any three of: calmness, flexibility, strength, balance, better body awareness, less worry while attention is on the practice.
\item (a) 3, (b) 1, (c) 4, (d) 2.
\item Yoga alone burns fewer calories than strong exercise. With sensible diet and some aerobic activity it can help control weight.
\item Any four of: empty or light bowel and bladder; a ventilated place; a level floor with a mat or blanket; unwanted thoughts kept aside; no practice during illness, after surgery, fracture or sprain until a physician allows.
\item Say it supports well-being: attention on the practice takes the mind off worries and pranayama brings calmness. Do not promise to reduce anxiety or cure.
\end{enumerate}

\paragraph{Sources.} The Board syllabus for Level 1 and the official guidebook chapter on Hatha Yoga (2016, older scheme, coverage check). The Hatha text of Svatmarama and the Gheranda Samhita, in a standard edition, are still to be used to confirm verse numbers.

\section{Mental well-being according to Patanjala yoga}
\label{sec:mental-wellbeing-patanjala}

Mental well-being in this view is a calm, clear mind that is not tossed by liking, dislike or fear. The four attitudes are the daily method. Learn the lists as named, and keep the Sanskrit term with one English word, because translations of the vrittis and the obstacle names differ between commentaries.

\subsection*{The mind and its waves}
Chitta is the mind-stuff. Its waves are vrittis, five in number: right knowledge, wrong knowledge, imagination, sleep, memory. Yoga is the stilling of these waves (chitta vritti nirodha).

\subsection*{Five states of the mind (chitta bhumi)}
Kshipta (scattered), mudha (dull), vikshipta (occasionally steady: mostly sattvic, but rajas pulls it away), ekagra (one-pointed), niruddha (stilled: all movements, even the sattvic ones, are held).

\subsection*{Two means of stilling the mind}
Abhyasa (steady practice) and vairagya (non-attachment). They are supported by social and personal discipline (yama, niyama), healthy diet, pranayama, asana and meditation.

\subsection*{Nine obstacles (antaraya, chitta vikshepa)}
\begin{itemize}
\item Disease.
\item Mental dullness.
\item Doubt.
\item Heedlessness.
\item Laziness.
\item Attachment to the senses.
\item Wrong understanding.
\item Failure to reach a firm ground.
\item Loss of steadiness once reached.
\end{itemize}

\subsection*{Four companions of the obstacles (sahabhuva)}
Unhappiness, despair, trembling of the body, irregular breathing.

\subsection*{Remedy named by the text}
Practice on one principle (eka-tattva abhyasa); repeating Om (pranava) with meaning; and the practice of the eight limbs.

\subsection*{Chitta prasadana: making the mind clear and calm}
Four attitudes:
\begin{itemize}
\item maitri: friendliness to the happy;
\item karuna: compassion to the suffering;
\item mudita: gladness for the virtuous;
\item upeksha: equanimity, a calm non-reaction towards the unvirtuous.
\end{itemize}
Write equanimity, not indifference, since indifference can be read as not caring. The breath way is also named: a long exhalation with a pause after it.

\subsection*{Five kleshas, the causes of distress}
Ignorance, ego-sense, liking, dislike, fear of loss (clinging to life). The guidebook says kriya yoga (tapas, svadhyaya, ishvara pranidhana) weakens them and meditation with pratiprasava removes them. Abhyasa and vairagya are the means for stilling the vrittis.

\subsection*{Sattva and pratipaksha bhavana}
\begin{itemize}
\item Sattva is the quality of purity, harmony and clarity of the mind. Yoga practice cultivates it.
\item Pratipaksha bhavana is answering a negative thought or feeling by deliberately holding its positive opposite.
\end{itemize}

\begin{caution}
Mindfulness of thoughts and feelings is a modern word for the same awareness. It is not a named practice in the guidebook. A qualified reviewer must check the health statements in this section.
\end{caution}

\begin{recall}
\begin{itemize}
\item Yoga is chitta vritti nirodha, the stilling of the waves of the mind.
\item Five vrittis, five states of mind, nine obstacles, four companions, five kleshas, four attitudes.
\item Two means: abhyasa and vairagya.
\item Four attitudes: maitri, karuna, mudita, upeksha (equanimity).
\item Sattva: purity, harmony, clarity.
\end{itemize}
\end{recall}

\subsection*{Common questions}
\begin{itemize}
\item \emph{Is upeksha indifference?} Use equanimity in exam answers.
\item \emph{Why do English names for the vrittis and obstacles differ between books?} Translations differ between commentaries. Keep the Sanskrit term with one English word.
\item \emph{Who elaborated the five states of the mind?} Vyasa, in his commentary on Yoga Sutra 1.1.
\end{itemize}

\subsection*{Worked examples}
\begin{itemize}
\item \textbf{Model answer: explain mental well-being according to Patanjala yoga.} Aim: a clear and steady mind; yoga is chitta vritti nirodha. Means: abhyasa and vairagya. Support: yama and niyama, healthy diet, pranayama, asana and meditation. Daily method: the four attitudes of chitta prasadana. Quality to grow: sattva, purity, harmony and clarity.
\item \textbf{Applying the four attitudes.} A friend is celebrating, a neighbour is grieving, a colleague is doing good work, and someone around you is behaving badly. Celebrating: maitri (friendliness). Grieving: karuna (compassion). Good work: mudita (gladness). Bad behaviour: upeksha (equanimity, a calm non-reaction).
\item \textbf{Memory aid for the lists.} Five states of mind, five vrittis, five kleshas, nine obstacles, four companions, four attitudes. Say each list with its count and its first and last item, then fill the middle.
\end{itemize}

\subsection*{Check yourself}
\begin{enumerate}
\item Name the two means of stilling the mind.
\item Name the five kleshas.
\item Fill in the blank: the five states of mind are kshipta, mudha, vikshipta, ekagra and \underline{\hspace{2cm}}.
\item Name the four companions of the obstacles.
\item What is pratipaksha bhavana?
\item Match the attitudes to the people: maitri, karuna, mudita, upeksha with (a) the suffering, (b) the happy, (c) the unvirtuous, (d) the virtuous.
\item True or false: sattva is the quality of restlessness.
\item A student feels low and finds breathing irregular and the body trembling. Which part of the Material does this match, and what remedy does the text name?
\item Which sutra idea defines yoga?
\end{enumerate}

\subsection*{Answers}
\begin{enumerate}
\item Abhyasa (practice) and vairagya (non-attachment).
\item Ignorance, ego-sense, liking, dislike and fear of loss (clinging to life).
\item Niruddha.
\item Unhappiness, despair, trembling of the body and irregular breathing.
\item Answering a negative thought or feeling by deliberately holding its positive opposite.
\item Maitri b, karuna a, mudita d, upeksha c.
\item False. Sattva is purity, harmony and clarity of the mind, and yoga practice cultivates it.
\item These are companions of the obstacles. The text names practice on one principle (eka-tattva abhyasa), repeating Om with meaning, and the practice of the eight limbs.
\item Chitta vritti nirodha, the stilling of the waves of the mind.
\end{enumerate}

\paragraph{Sources.} The Board syllabus for Level 1 and the official guidebook chapters on Patanjala Yoga and stress management (2016, older scheme, coverage check). A standard edition of the Yoga Sutra with commentary is still to be set against the sutra names and numbers.

\section{Bahiranga and antaranga yoga}
\label{sec:bahiranga-antaranga}

These are the outer and inner limbs of the eightfold path (ashtanga). Name all eight in order and say which are outer and which are inner.

\subsection*{The eight limbs}
Yama (restraints), niyama (observances), asana (posture), pranayama (breath control), pratyahara (withdrawal of the senses), dharana (holding attention), dhyana (meditation), samadhi (absorption).

\subsection*{Outer and inner}
\begin{itemize}
\item Outer (bahiranga): the first five, yama, niyama, asana, pranayama, pratyahara.
\item Inner (antaranga): the last three, dharana, dhyana, samadhi, counted together as samyama in the third pada.
\end{itemize}
The Level 1 syllabus gives no split, so teach five outer and three inner. Some books count only four outer limbs and call pratyahara a bridge between the first four and the last three. You may mention that view.

\subsection*{Yama and niyama in a line each}
\begin{itemize}
\item Yamas: non-harming, truthfulness, not stealing, moderation of the senses, non-possessiveness.
\item Niyamas: cleanliness, contentment, discipline, self-study, surrender to the ideal.
\end{itemize}

\subsection*{One-line role of each limb}
\begin{itemize}
\item Yama: ethical (social) restraints.
\item Niyama: personal observances.
\item Asana: a posture that prepares the body for meditation.
\item Pranayama: breath control that steadies energy and the mind.
\item Pratyahara: withdrawal of the senses.
\item Dharana: attention held on one point.
\item Dhyana: unbroken flow of attention, meditation.
\item Samadhi: absorption, deep inner stillness. Write absorption, not enlightenment.
\end{itemize}

\subsection*{How the groups relate}
Use the ladder picture: each limb prepares the next. Say the outer limbs prepare the inner. Steady practice of the limbs is said to thin impurities until discriminative knowledge dawns.

\begin{recall}
\begin{itemize}
\item Order: yama, niyama, asana, pranayama, pratyahara, dharana, dhyana, samadhi.
\item Outer: the first five. Inner: dharana, dhyana, samadhi, together samyama.
\item Samadhi is absorption.
\item Ladder: outer ones prepare the inner.
\end{itemize}
\end{recall}

\subsection*{Common questions}
\begin{itemize}
\item \emph{What is samyama?} The name for dharana, dhyana and samadhi taken together.
\item \emph{Do I have to learn the limbs in order?} Yes. Each limb prepares the next.
\item \emph{Are the outer and inner limbs separate paths?} No. Put it as the outer limbs preparing the inner.
\end{itemize}

\subsection*{Worked examples}
\begin{itemize}
\item \textbf{Model answer: name the eight limbs and say which are outer and which are inner.} Write the eight in order, then draw the line after the fifth. Outer: yama to pratyahara. Inner: dharana, dhyana, samadhi, together samyama in the third pada. Add one line: the outer limbs prepare the inner.
\item \textbf{Memory aid: the ladder.} Say the outer steps as ethics, ethics, body, breath, senses. Then the inner steps as hold, flow, absorb.
\item \textbf{The pratyahara question.} A book counts four outer limbs and calls pratyahara a bridge. Remember the Level 1 syllabus gives no split. In the exam write five outer and three inner, which matches the Sutra. You may add the bridge view.
\end{itemize}

\subsection*{Check yourself}
\begin{enumerate}
\item Fill in the blank: the eight limbs are yama, niyama, asana, pranayama, pratyahara, dharana, dhyana and \underline{\hspace{2cm}}.
\item Name the five yamas.
\item Name the five niyamas.
\item True or false: samadhi means enlightenment.
\item Match each to its one-line role: (a) asana, (b) pranayama, (c) dharana, (d) dhyana.
\item What is the difference between yama and niyama?
\item Why are both groups of limbs needed?
\item A student asks which limbs are inner and what they are called together. How do you answer?
\end{enumerate}

\subsection*{Answers}
\begin{enumerate}
\item Samadhi.
\item Non-harming, truthfulness, not stealing, moderation of the senses, non-possessiveness.
\item Cleanliness, contentment, discipline, self-study and surrender to the ideal.
\item False as the exam wording. Write absorption (deep inner stillness) instead.
\item (a) a posture that prepares the body for meditation, (b) breath control that steadies energy and the mind, (c) attention held on one point, (d) unbroken flow of attention, meditation.
\item Yama is ethical (social) restraint; niyama is personal observance.
\item They support each other; the outer limbs prepare the ground for the inner.
\item Dharana, dhyana and samadhi; together they are samyama (third pada).
\end{enumerate}

\paragraph{Sources.} The Board syllabus for Level 1 and the official guidebook chapter on Patanjala Yoga (2016, older scheme, coverage check). The Board syllabi for the higher levels give the same five and three split.
